\documentclass[11pt]{article}
\usepackage[margin=0.9in]{geometry}
\usepackage{listings}
\usepackage{xcolor}
\usepackage{amssymb}
\usepackage{enumitem}
\usepackage{fancyhdr}
\usepackage{titlesec}

% ---------- Colors ----------
\definecolor{codebg}{RGB}{245,245,245}
\definecolor{codekeyword}{RGB}{0,90,150}
\definecolor{codestring}{RGB}{140,20,20}
\definecolor{codecomment}{RGB}{100,100,100}

% ---------- Code listing style (Python) ----------
\lstdefinestyle{pystyle}{
    language=Python,
    backgroundcolor=\color{codebg},
    basicstyle=\ttfamily\small,
    keywordstyle=\color{codekeyword}\bfseries,
    stringstyle=\color{codestring},
    commentstyle=\color{codecomment}\itshape,
    showstringspaces=false,
    breaklines=true,
    frame=single,
    rulecolor=\color{gray!40},
    numbers=left,
    numberstyle=\tiny\color{gray},
    xleftmargin=1.8em,
    framexleftmargin=1.5em,
    tabsize=4,
}
\lstset{style=pystyle}

% ---------- Header/Footer ----------
\pagestyle{fancy}
\setlength{\headheight}{14pt}
\fancyhf{}
\lhead{Control Flow: Predict the Output}
\rhead{Name: \underline{\hspace{4cm}}}
\cfoot{\thepage}

\titleformat{\section}{\large\bfseries}{\thesection}{0.5em}{}
\titleformat{\subsection}{\normalsize\bfseries}{\thesubsection}{0.5em}{}

\newcommand{\answerbox}{%
\vspace{2mm}
\noindent\textbf{Output:}
\vspace{14mm}
}

\begin{document}

\begin{center}
    {\Large \textbf{Control Flow Detective}}\\[2mm]
    {\normalsize Read each piece of code carefully. Do NOT run it in your head too fast --- trace through it
    line by line, just like the computer would. Write down exactly what would be printed to the console.}
\end{center}

\vspace{4mm}
\noindent\textbf{Directions:} For each problem, write the exact output, including spacing, capitalization,
and line breaks. If nothing is printed, write \texttt{(nothing)}. If the code would cause an error,
write \texttt{ERROR} and explain why in one sentence.

\vspace{4mm}
\hrule
\vspace{4mm}

\noindent\textbf{Worked Example} \quad (this one is done for you --- study how the tracing works!)

\begin{lstlisting}
x = 3
if x > 5:
    print("big")
else:
    print("small")
print("done")
\end{lstlisting}

\noindent\textit{How to trace it:}
\begin{enumerate}[nosep]
    \item \texttt{x} is set to \texttt{3}.
    \item Check the condition: is \texttt{3 > 5}? No, that's false.
    \item Since the \texttt{if} was false, we jump to the \texttt{else} block and run \texttt{print("small")}.
    \item The \texttt{if}/\texttt{else} is finished, so we move to the next line, \texttt{print("done")}, which always runs no matter which branch we took.
\end{enumerate}

\noindent\textbf{Output:}
\begin{lstlisting}[frame=none,numbers=none,backgroundcolor=\color{white}]
small
done
\end{lstlisting}

\vspace{4mm}
\hrule

\vspace{2mm}

\vspace{4mm}
\hrule
\vspace{2mm}

\noindent\textbf{Quick reminders:}
\vspace{2mm}

\begin{itemize}[nosep]
    \item \texttt{\%} is the \textbf{modulo} operator --- it gives you the \textit{remainder} after
    dividing. For example, \texttt{7 \% 2} is \texttt{1} (7 divided by 2 is 3 with 1 left over), and
    \texttt{10 \% 5} is \texttt{0} (10 divides evenly by 5, so there's no remainder). A number is
    \textbf{even} if \texttt{n \% 2 == 0}, and \textbf{odd} if \texttt{n \% 2 == 1}.
    \item \texttt{**} is the \textbf{exponent} operator --- it means "raised to the power of."
    For example, \texttt{2 ** 3} means $2^3$, which is \texttt{8}, and \texttt{5 ** 2} means
    $5^2$, which is \texttt{25}.
\end{itemize}

\vspace{2mm}
\hrule
\vspace{2mm}

% =========================================================
\section{Conditionals}

\vspace{4mm}
\noindent\textbf{Problem 1}
\begin{lstlisting}
x = 7
if x > 5:
    print("big")
else:
    print("small")
\end{lstlisting}
\answerbox

\vspace{4mm}
\noindent\textbf{Problem 2}
\begin{lstlisting}
temperature = 68

if temperature > 80:
    print("It's hot!")
elif temperature > 60:
    print("Nice day.")
else:
    print("Bring a jacket.")
\end{lstlisting}
\answerbox

\vspace{4mm}
\noindent\textbf{Problem 3}
\begin{lstlisting}
grade = 72

if grade >= 90:
    print("A")
elif grade >= 80:
    print("B")
elif grade >= 70:
    print("C")
else:
    print("F")
\end{lstlisting}
\answerbox

\vspace{4mm}
\noindent\textbf{Problem 4 (tricky!)}
\begin{lstlisting}
password = "wizard123"

if password == "wizard123":
    print("Access granted")
if password == "Wizard123":
    print("Access granted (capital W)")
else:
    print("Access denied (capital W)")
\end{lstlisting}
\answerbox

% =========================================================
\newpage
\section{Loops}

\noindent\textbf{A quick note on \texttt{break} and \texttt{continue}:}
\vspace{2mm}

\noindent Sometimes you want to change how a loop behaves \textit{while it's running}, instead of
just letting it finish normally. Two keywords let you do that:

\begin{itemize}[nosep]
    \item \texttt{break} --- immediately \textbf{stops the loop completely}. The program jumps to
    whatever comes after the loop, even if there were more values left to go through.
    \item \texttt{continue} --- \textbf{skips the rest of the current pass} through the loop and
    jumps straight to the next one. It does NOT stop the loop --- it just skips ahead.
\end{itemize}

\begin{lstlisting}
for i in range(5):
    if i == 3:
        break
    print(i)
\end{lstlisting}
\noindent This prints \texttt{0 1 2} and then stops --- once \texttt{i} hits \texttt{3}, the loop
ends completely.

\begin{lstlisting}
for i in range(5):
    if i == 3:
        continue
    print(i)
\end{lstlisting}
\noindent This prints \texttt{0 1 2 4} --- when \texttt{i} is \texttt{3}, that one pass is
skipped, but the loop keeps going afterward.

\vspace{4mm}
\hrule
\vspace{2mm}

\vspace{4mm}
\noindent\textbf{Problem 5}
\begin{lstlisting}
for i in range(5):
    print(i)
\end{lstlisting}
\answerbox

\vspace{4mm}
\noindent\textbf{Problem 6}
\begin{lstlisting}
count = 0
while count < 3:
    print("Hello!")
    count = count + 1
\end{lstlisting}
\answerbox

\vspace{4mm}
\noindent\textbf{Problem 7}
\begin{lstlisting}
total = 0
for num in range(1, 5):
    total = total + num
print(total)
\end{lstlisting}
\answerbox

\vspace{4mm}
\noindent\textbf{Problem 8 (break and continue)}
\begin{lstlisting}
for i in range(10):
    if i == 3:
        continue
    if i == 6:
        break
    print(i)
\end{lstlisting}
\answerbox

\vspace{4mm}
\noindent\textbf{Problem 9 (nested loop)}
\begin{lstlisting}
for i in range(3):
    for j in range(2):
        print(i, j)
\end{lstlisting}
\answerbox

% =========================================================
\newpage
\section{Functions}

\vspace{4mm}
\noindent\textbf{Problem 10}
\begin{lstlisting}
def greet(name):
    print("Hi, " + name + "!")

greet("Maryam")
greet("Amina")
\end{lstlisting}
\answerbox

\vspace{4mm}
\noindent\textbf{Problem 11}
\begin{lstlisting}
def square(x):
    return x * x

result = square(4)
print(result)
\end{lstlisting}
\answerbox

\vspace{4mm}
\noindent\textbf{Problem 12 (return vs.\ print --- read carefully!)}
\begin{lstlisting}
def mystery(n):
    if n % 2 == 0:
        return "even"
    else:
        return "odd"

print(mystery(7))
print(mystery(10))
\end{lstlisting}
\answerbox

\vspace{4mm}
\noindent\textbf{Problem 13 (default arguments)}
\begin{lstlisting}
def power(base, exponent=2):
    return base ** exponent

print(power(3))
print(power(2, 3))
\end{lstlisting}
\answerbox

% =========================================================
\newpage
\section{Combining Everything}

\vspace{4mm}
\noindent\textbf{Problem 14}
\begin{lstlisting}
def check_number(n):
    if n < 0:
        return "negative"
    elif n == 0:
        return "zero"
    else:
        return "positive"

numbers = [-3, 0, 8, -1]
for num in numbers:
    print(check_number(num))
\end{lstlisting}
\answerbox

\vspace{4mm}
\noindent\textbf{Problem 15}
\begin{lstlisting}
def countdown(n):
    while n > 0:
        print(n)
        n = n - 1
    print("Liftoff!")

countdown(3)
\end{lstlisting}
\answerbox

\vspace{4mm}
\noindent\textbf{Problem 16 (challenge: functions calling functions)}
\begin{lstlisting}
def double(x):
    return x * 2

def triple(x):
    return x * 3

def combo(x):
    if x % 2 == 0:
        return double(x)
    else:
        return triple(x)

for i in range(1, 5):
    print(combo(i))
\end{lstlisting}
\answerbox

\vspace{4mm}
\noindent\textbf{Problem 17 (challenge: nested loop + conditional)}
\begin{lstlisting}
for i in range(1, 4):
    for j in range(1, 4):
        if (i + j) % 2 == 0:
            print(str(i) + "+" + str(j) + " = even sum")
\end{lstlisting}
\answerbox

\vspace{4mm}
\noindent\textbf{Problem 18 (challenge: FizzBuzz)}
\begin{lstlisting}
for n in range(1, 16):
    if n % 15 == 0:
        print("FizzBuzz")
    elif n % 3 == 0:
        print("Fizz")
    elif n % 5 == 0:
        print("Buzz")
    else:
        print(n)
\end{lstlisting}
\answerbox

\end{document}
